\subsection{Quotient Modules, Correspondence, and Isomorphism}

Throughout this subsection, $R$ is a nonzero commutative ring with identity $1_R$.\\
Fix an $R$-module $M$ with additive identity $0_M$.

\begin{definition} \label{def:quotient-module} \hlt{Quotient Module}\\
Let $S\leq M$. For $(u,v)\in M^2$, define $u\equiv v\pmod S\Longleftrightarrow u-v\in S$.\\
For $u\in M$, set $u+S\defeq\{u+s:s\in S\}$, and set $M/S\defeq\{u+S:u\in M\}$.\\
For $(u,v)\in M^2$, set $(u+S)+(v+S)\defeq(u+v)+S$. For $(r,u)\in R\times M$, set $r(u+S)\defeq ru+S$.
\end{definition}

\begin{proposition} \label{prop:quotient-module-well-defined} \hlt{Quotient Operations Are Well-Defined}\\
Let $S\leq M$.
\begin{enumerate}[label=(\roman*)]
\setlength{\itemsep}{0pt}
\item For $(u,v)\in M^2$, $u+S=v+S\Longleftrightarrow u-v\in S$. Consequently, congruence modulo $S$ is an equivalence relation.
\item The operations of Definition \ref{def:quotient-module} are well-defined and make $M/S$ an $R$-module.
\item The additive identity is $0_{M/S}=S$. Quotient map $\pi:M\twoheadrightarrow M/S$, $\pi(u)\defeq u+S$, is epimorphism with $\ker\pi=S$.
\end{enumerate}
\begin{proof}
\begin{enumerate}[label=(\roman*)]
\setlength{\itemsep}{0pt}
\item[$(i)$] ($\Rightarrow$) Assume $u+S=v+S$. Then $u\in u+S=v+S$ by $0_M\in S$. Thus $(\exists\,s\in S)\,u=v+s$ by Definition \ref{def:quotient-module}, hence $u-v\in S$.
\item[$(i)$] ($\Leftarrow$) Assume $u-v\in S$ and fix $x\in u+S$. Choose $s\in S$ with $x=u+s$ by Definition \ref{def:quotient-module}.\\
$x=v+\bigl((u-v)+s\bigr)\in v+S$ by closure of $S$, so $u+S\subseteq v+S$. Fix $y\in v+S$, choose $t\in S$ with $y=v+t$ by Definition \ref{def:quotient-module}.\\
Then $y=u+\bigl((v-u)+t\bigr)\in u+S$ by $v-u=-(u-v)\in S$ and closure of $S$. Thus $v+S\subseteq u+S$ and $u+S=v+S$.\\
Congruence modulo $S$ is an equivalence relation by this characterization and equality of cosets.
\item[$(ii)$] Assume $u+S=u'+S$ and $v+S=v'+S$. Then $u-u',v-v'\in S$ by (i). $(u+v)-(u'+v')=(u-u')+(v-v')\in S$.\\
For $r\in R$, $ru-ru'=r(u-u')\in S$. Hence $(u+v)+S=(u'+v')+S$ and $ru+S=ru'+S$ by (i).\\
Thus both operations are representative-independent.\\
Every module axiom holds on representatives by the corresponding axiom in $M$. Therefore $M/S$ is an $R$-module.
\item[$(iii)$] The additive identity is $0_{M/S}=S$ by $0_M+S=S$.\\
For $(r,s,u,v)\in R^2\times M^2$, $\pi(ru+sv)=(ru+sv)+S=r(u+S)+s(v+S)=r\pi(u)+s\pi(v)$.\\
Thus $\pi\in\operatorname{Hom}_R(M,M/S)$ by Definition \ref{def:module-homomorphism}(i). The map $\pi$ is surjective by $M/S=\{u+S:u\in M\}$.\\
Hence $\pi$ is an epimorphism by Definition \ref{def:module-homomorphism}(iv). Finally, $\ker\pi=\{u\in M:u+S=S\}=\{u\in M:u\in S\}=S$ by (i).
\end{enumerate}
\hlt{Structural Commentary}: Coset equality controls representatives; module axioms descend from $M$.\\
\hlt{Key Insight}: Representatives agree exactly when their difference lies in $S$.
\end{proof}
\end{proposition}

\begin{theorem} \label{thm:module-correspondence-theorem} \hlt{Correspondence Theorem for Modules}\\
Let $S\leq M$ and let $\pi:M\twoheadrightarrow M/S$ be the quotient map. Set $\mathcal{A}\defeq\{T\leq M:S\leq T\}$ and $\mathcal{B}\defeq\{U:U\leq M/S\}$.\\
For $T\in\mathcal{A}$, define $\Phi(T)\defeq T/S=\pi(T)$. For $U\in\mathcal{B}$, define $\Psi(U)\defeq\pi^{-1}(U)$.
\begin{enumerate}[label=(\roman*)]
\setlength{\itemsep}{0pt}
\item $(\forall\,T\in\mathcal{A})\,\Phi(T)\in\mathcal{B}$ and $(\forall\,U\in\mathcal{B})\,\Psi(U)\in\mathcal{A}$.
\item $\Psi\Phi=\operatorname{id}_{\mathcal{A}}$ and $\Phi\Psi=\operatorname{id}_{\mathcal{B}}$.
\item Both $\Phi$ and $\Psi$ preserve inclusion.
\end{enumerate}
Consequently, $\Phi:\mathcal{A}\to\mathcal{B}$ is an order isomorphism with inverse $\Psi$.
\begin{proof}
\begin{enumerate}[label=(\roman*)]
\setlength{\itemsep}{0pt}
\item Fix $T\in\mathcal{A}$; $0_{M/S}=S\in T/S$ by $S\leq T$. For $(r,s,t_1,t_2)\in R^2\times T^2$, $r(t_1+S)+s(t_2+S)=(rt_1+st_2)+S$ by Definition \ref{def:quotient-module}.\\
The coset $(rt_1+st_2)+S$ lies in $T/S$ by closure of $T$. Thus $\Phi(T)=T/S\leq M/S$ by Theorem \ref{thm:submodule-test}. Fix $U\in\mathcal{B}$.\\
Then $0_M\in\pi^{-1}(U)$ by $\pi(0_M)=0_{M/S}\in U$. Fix $(r,s,u,v)\in R^2\times\pi^{-1}(U)^2$.\\
Then $\pi(ru+sv)=r\pi(u)+s\pi(v)\in U$ by Proposition \ref{prop:quotient-module-well-defined}(iii) and closure of $U$. Hence $\pi^{-1}(U)\leq M$ by Theorem \ref{thm:submodule-test}.\\
$S=\ker\pi\subseteq\pi^{-1}(U)$ by Proposition \ref{prop:quotient-module-well-defined}(iii) and $0_{M/S}\in U$. Hence $\Psi(U)=\pi^{-1}(U)\in\mathcal{A}$ by the definition of $\mathcal{A}$.
\item Fix $T\in\mathcal{A}$. Then $(\forall\,t\in T)\,\pi(t)=t+S$ by Proposition \ref{prop:quotient-module-well-defined}(iii). Hence $T\subseteq\pi^{-1}(T/S)$ by the definition of $T/S$.\\
If $m\in\pi^{-1}(T/S)$, choose $t\in T$ with $m+S=t+S$ by the definition of $T/S$.\\
Then $m-t\in S\subseteq T$ by Proposition \ref{prop:quotient-module-well-defined}(i), so $m\in T$. Therefore $\Psi\Phi(T)=\pi^{-1}(T/S)=T$. Fix $U\in\mathcal{B}$.\\
Inclusion $\pi(\pi^{-1}(U))\subseteq U$ holds by definition of preimage. If $x\in U$, choose $m\in M$ with $x=\pi(m)$ by Proposition \ref{prop:quotient-module-well-defined}(iii).\\
Then $m\in\pi^{-1}(U)$ and $x\in\pi(\pi^{-1}(U))$. Therefore $\Phi\Psi(U)=\pi(\pi^{-1}(U))=U$.
\item If $T_1\subseteq T_2$ in $\mathcal{A}$, then $\Phi(T_1)=\pi(T_1)\subseteq\pi(T_2)=\Phi(T_2)$. If $U_1\subseteq U_2$ in $\mathcal{B}$, then $\Psi(U_1)=\pi^{-1}(U_1)\subseteq\pi^{-1}(U_2)=\Psi(U_2)$.\\
Thus both assignments preserve inclusion.
\end{enumerate}
Thus $\Phi:\mathcal{A}\to\mathcal{B}$ is an order isomorphism with inverse $\Psi$ by (ii)--(iii).\\
\hlt{Structural Commentary}: Closure constructs $\Phi,\Psi$; lifting identifies inverses; set maps preserve order.\\
\hlt{Key Insight}: Preimages recover quotient images by surjectivity and kernel containment.
\end{proof}
\end{theorem}

\begin{theorem} \label{thm:module-isomorphism-theorems} \hlt{Module Isomorphism Theorems}\\
The following canonical maps are $R$-module isomorphisms.
\begin{enumerate}[label=(\roman*)]
\setlength{\itemsep}{0pt}
\item Let $N$ be $R$-module with additive identity $0_N$, $\tau\in\operatorname{Hom}_R(M,N)$. For $m\in M$, $\overline{\tau}(m+\ker\tau)\defeq\tau(m)$. Then $\overline{\tau}:M/\ker\tau\xrightarrow{\sim}\operatorname{im}\tau$.
\item For $S,T\leq M$ and $s\in S$, define $\theta\bigl(s+(S\cap T)\bigr)\defeq s+T$. Then $\theta:S/(S\cap T)\xrightarrow{\sim}(S+T)/T$.
\item For $S\leq T\leq M$ and $m\in M$, define $\eta\bigl((m+S)+T/S\bigr)\defeq m+T$. Then $\eta:(M/S)/(T/S)\xrightarrow{\sim}M/T$.
\end{enumerate}
\begin{proof}
\begin{enumerate}[label=(\roman*)]
\setlength{\itemsep}{0pt}
\item Assume $m+\ker\tau=m'+\ker\tau$. Then $m-m'\in\ker\tau$ by Proposition \ref{prop:quotient-module-well-defined}(i).\\
$\tau(m)-\tau(m')=\tau(m-m')=0_N$ by Definition \ref{def:module-homomorphism}(i). Thus $\tau(m)=\tau(m')$ and $\overline{\tau}$ is well-defined.\\
Fix $(r,s,m,m')\in R^2\times M^2$. Then $\overline{\tau}\bigl(r(m+\ker\tau)+s(m'+\ker\tau)\bigr)=\tau(rm+sm')=r\tau(m)+s\tau(m')$.\\
Therefore $\overline{\tau}\in\operatorname{Hom}_R(M/\ker\tau,\operatorname{im}\tau)$ by Definition \ref{def:module-homomorphism}(i). $\operatorname{im}\overline{\tau}=\operatorname{im}\tau$ by the definition of $\overline{\tau}$.\\
$\ker\overline{\tau}=\{m+\ker\tau:m\in\ker\tau\}=\{\ker\tau\}=\{0_{M/\ker\tau}\}$. Thus $\overline{\tau}$ is a monomorphism by Proposition \ref{prop:hom-module-kernel-image}(iv).\\
It is an epimorphism by $\operatorname{im}\overline{\tau}=\operatorname{im}\tau$ and Definition \ref{def:module-homomorphism}(iv). Hence $\overline{\tau}$ is an isomorphism by Definition \ref{def:module-homomorphism}(iv).
\item Assume $s+(S\cap T)=s'+(S\cap T)$. Then $s-s'\in S\cap T\subseteq T$ by Proposition \ref{prop:quotient-module-well-defined}(i).\\
Thus $s+T=s'+T$ by Proposition \ref{prop:quotient-module-well-defined}(i), so $\theta$ is well-defined. Fix $(r,q,s_1,s_2)\in R^2\times S^2$.\\
Then $\theta\bigl(r(s_1+(S\cap T))+q(s_2+(S\cap T))\bigr)=r(s_1+T)+q(s_2+T)$. Hence $\theta\in\operatorname{Hom}_R(S/(S\cap T),(S+T)/T)$ by Definition \ref{def:module-homomorphism}(i).\\
Fix $x\in(S+T)/T$, $y\in S+T$ with $x=y+T$ by Definition \ref{def:quotient-module}. Choose $(s,t)\in S\times T$ with $y=s+t$ by Proposition \ref{prop:submodule-intersection-sum}(ii).\\
Then $x=s+T=\theta\bigl(s+(S\cap T)\bigr)$, so $\theta$ is surjective. $\ker\theta=\{s+(S\cap T):s\in S\cap T\}=\{S\cap T\}$ by Proposition \ref{prop:quotient-module-well-defined}(i).\\
Hence $\ker\theta=\{0_{S/(S\cap T)}\}$ by $0_{S/(S\cap T)}=S\cap T$. Thus $\theta$ is an isomorphism by Proposition \ref{prop:hom-module-kernel-image}(iv) and Definition \ref{def:module-homomorphism}(iv).
\item Assume $(m+S)+T/S=(m'+S)+T/S$. Then $(m-m')+S\in T/S$ by Proposition \ref{prop:quotient-module-well-defined}(i).\\
Choose $t\in T$ with $(m-m')+S=t+S$ by the definition of $T/S$. $m-m'-t\in S\subseteq T$ by Proposition \ref{prop:quotient-module-well-defined}(i), so $m-m'\in T$.\\
Thus $m+T=m'+T$ by Proposition \ref{prop:quotient-module-well-defined}(i), and $\eta$ is well-defined. Fix $(r,s,m,m')\in R^2\times M^2$.\\
Then $\eta\bigl(r((m+S)+T/S)+s((m'+S)+T/S)\bigr)=r(m+T)+s(m'+T)$; $\eta\in\operatorname{Hom}_R((M/S)/(T/S),M/T)$ by Definition \ref{def:module-homomorphism}(i).\\
$\eta$ is surjective by $(\forall\,m\in M)\,\eta\bigl((m+S)+T/S\bigr)=m+T$. $\ker\eta=\{(m+S)+T/S:m\in T\}=\{T/S\}$ by Proposition \ref{prop:quotient-module-well-defined}(i).\\
Hence $\ker\eta=\{0_{(M/S)/(T/S)}\}$ by $0_{(M/S)/(T/S)}=T/S$. Thus $\eta$ is isomorphism by Proposition \ref{prop:hom-module-kernel-image}(iv), Definition \ref{def:module-homomorphism}(iv).
\end{enumerate}
\hlt{Structural Commentary}: Coset checks define maps; kernels and representatives certify bijectivity.\\
\hlt{Key Insight}: Each denominator is exactly the ambiguity annihilated by its canonical map.
\end{proof}
\end{theorem}
